\documentclass[12pt,a4paper]{article}

\usepackage[utf8]{inputenc}
\usepackage[T1]{fontenc}
\usepackage[ngerman]{babel}
\usepackage{lmodern}
\usepackage[a4paper,left=2.8cm,right=2.3cm,top=2.5cm,bottom=2.5cm]{geometry}
\usepackage{microtype}
\usepackage{graphicx}
\usepackage{booktabs}
\usepackage{tabularx}
\usepackage{longtable}
\usepackage{array}
\usepackage{enumitem}
\usepackage{amsmath}
\usepackage{xcolor}
\usepackage{float}
\usepackage{pdflscape}
\usepackage{fancyhdr}
\usepackage{xurl}
\usepackage{tikz}
\usetikzlibrary{arrows.meta,positioning,fit,backgrounds,shapes.geometric}
\let\tikzpathcommand\path
\renewcommand{\path}[1]{\nolinkurl{#1}}
\usepackage[breaklinks=true]{hyperref}
\usepackage{csquotes}
\usepackage[backend=biber,style=authoryear,sorting=nyt]{biblatex}

\newcommand{\thwsAuthor}{[Martin Rapps]}
\newcommand{\thwsMatriculation}{[6323014]}
\newcommand{\thwsDegree}{[B. Eng. Geovisualisierung]}
\newcommand{\thwsFirstExaminer}{[Dr. Markus Müller]}
\newcommand{\thwsSecondExaminer}{[Andreas Rupp]}
\newcommand{\thwsSubmissionDate}{[08.09.2026]}

\hypersetup{colorlinks=true,linkcolor=blue,citecolor=blue,urlcolor=blue}
\setlist{itemsep=2pt,topsep=4pt}
\pagestyle{fancy}
\fancyhf{}
\fancyhead[L]{Anlage -- Dokumentation der KI-Nutzung}
\fancyhead[R]{Scan-to-BIM-Pipeline}
\fancyfoot[C]{\thepage}
\setlength{\headheight}{15pt}

\begin{document}
\pagenumbering{Roman}
\thispagestyle{empty}
\begin{center}
	{\large Technische Hochschule Würzburg-Schweinfurt\\
	Fakultät Kunststofftechnik und Vermessung\par}
	\vfill
	{\LARGE\bfseries Anlage zur Projektarbeit\par}
	\vspace{1.0cm}
	{\Large\bfseries Dokumentation der KI-Nutzung\par}
	\vspace{0.6cm}
	{\large Scan-to-BIM-Pipeline -- KI-gestützte 3D-Rekonstruktion eines linearen Objekts\par}
	\vfill
	\begin{tabular}{@{}ll@{}}
		Vorgelegt von: & \thwsAuthor\\
		Matrikelnummer: & \thwsMatriculation\\
		Studiengang: & \thwsDegree\\
		Erstprüfer/in: & \thwsFirstExaminer\\
		Zweitprüfer/in: & \thwsSecondExaminer\\
		Abgabetermin: & \thwsSubmissionDate
	\end{tabular}
	\vfill
	Würzburg, den \today
\end{center}
\clearpage

\section*{Hinweis zur Dokumentationsform}
Auf Grund des Umfangs und der agentischen Arbeitsweise (selbstständige Änderungen über mehrere Dateien, iterative Überarbeitungen) wird die KI-Nutzung nach Bereichen beschrieben und jeweils durch beispielhafte Prompts veranschaulicht. Eine vollständige Auflistung aller Einzel-Prompts ist bei diesem Arbeitsmodus nicht sinnvoll und wurde so mit dem Prüfer abgestimmt.

Die Verantwortung für fachliche Richtigkeit, Eigenständigkeit und kritische Prüfung aller Inhalte liegt beim Verfasser. KI-Ausgaben wurden nicht als wissenschaftliche Quellen zitiert.

\tableofcontents
\clearpage
\pagenumbering{arabic}

\section{Projektangaben}

\begin{tabular}{@{}ll@{}}
Projekt: & Scan-to-BIM-Pipeline -- KI-gestützte 3D-Rekonstruktion eines linearen Objekts \\
Verfasser: & \thwsAuthor\ (\thwsMatriculation) \\
Zeitraum der KI-Nutzung: & Juni -- August 2026 \\
Abgabe: & \thwsSubmissionDate \\
\end{tabular}

\section{Eingesetzte Werkzeuge}

\begin{table}[H]
\centering
\small
\begin{tabularx}{\textwidth}{p{3.8cm}p{4.2cm}X}
\toprule
Werkzeug & Version / Anbieter & Einsatzbereich\\
\midrule
\texttt{opencode} (CLI-Agent) & aktuelle Version, verschiedene LLMs je nach Aufgabe & Code-Erstellung, Refactoring, Pipeline-Debugging, Versuchsplanung, Text- und Grafikarbeit\\
VS Code & aktuelle Version, integrierte KI-Unterstützung & Code-Editierung, Inline-Vorschläge, LaTeX-Bearbeitung\\
\bottomrule
\end{tabularx}
\caption{Übersicht der eingesetzten KI-Werkzeuge.}
\end{table}

Die Modellwahl innerhalb von \texttt{opencode} erfolgte je nach Aufgabenkomplexität und Kosten -- für einfache Korrekturen schlanke Modelle, für konzeptionelle oder fachliche Fragen leistungsfähigere Modelle.

\textbf{Beispielhafte Modelle:} Hauptsächlich genutzt wurden GPT-5.6 Luna (max), Gemini 3.5 Flash (max), Gemini 3.7 Flash (max), Kimi K3 (max), GLM 5.3 Flash (max) und GPT-5.6 Sol (max).

Alle KI-Ausgaben wurden manuell geprüft, verstanden und verantwortet.

\section{Code-Erstellung / Verbesserung / Refactoring}

\textbf{Worum es ging:} Aufbau und Pflege der Docker-basierten Pipeline (SAM-3.1-Segmentierung, COLMAP, STS, SuGaR-Fork, Centerline-Extraktion), insbesondere die Anpassung der Maskenlogik (Dilatation/Erosion für genau ein Zielobjekt), Fehlerbehandlung in Matrixläufen und die maskenbewusste Erweiterung des SuGaR-Forks.

\textbf{Beispielhafte Prompts:}
\begin{itemize}
  \item \enquote{Die SAM-Masken sollen als Hierarchie vorliegen: \texttt{default} unverändert, \texttt{middle} einmal mit einem $5\times5$-Fenster erodiert, \texttt{small} zweifach erodiert, ohne zusätzliche Dilatation. Stelle zusammen mit der promptbasierten Auswahl sicher, dass automatisch genau ein Zielobjekt übrig bleibt. Zeige nur die geänderte Funktion und erkläre die Parameter.}
  \item \enquote{Der Matrixrunner führt von meiner TSV-Konfiguration nur die erste Variante aus und beendet sich danach. Finde die Ursache und schlage eine Korrektur vor, die alle geplanten Läufe nacheinander ausführt.}
  \item \enquote{Im SuGaR-Fork soll der Verlust nur innerhalb der Objektmaske gewertet werden. Implementiere einen maskierten Loss mit Normierung über die Maskenpixel und Schutz vor Division durch null.}
\end{itemize}

\textbf{Prüfung:} Jeder Vorschlag wurde im lokalen Docker-Setup ausgeführt und gegen Manifeste/Logs geprüft; unverstandene Änderungen wurden nicht übernommen.

\section{Allgemeine Überlegungen zum Vorgehen (Sparringspartner)}

\textbf{Worum es ging:} Struktur der Arbeit, Versuchsplanung (Matrixdesign Kameramodell $\times$ FPS $\times$ Auflösung), Bewertung von Varianten (Route A vs.\ SuGaR-Coarse, Tiefenrouten), Interpretation von Metriken.

\textbf{Beispielhafte Prompts:}
\begin{itemize}
  \item \enquote{Ich habe die Kameramodelle \texttt{PINHOLE} und \texttt{OPENCV} getestet, beide liefern ähnliche Metriken. Erkläre fachlich, welche Argumente für welchen Standard sprechen und wo die Grenzen dieser Aussage ohne unabhängige Kalibrierung liegen.}
  \item \enquote{Hilf mir, die Matrixläufe so zu planen, dass ich Kameramodell, FPS und Auflösung getrennt auswerten kann, ohne die Meshroute zu vermischen.}
\end{itemize}

\textbf{Prüfung:} Vorschläge wurden mit dem Betreuer und dem Zielrahmen abgeglichen; fachliche Entscheidungen (z.\,B.\ Produktionsstandard) blieben eigene Arbeit.

\section{Schreiben und Überarbeiten von Texten in der Projektarbeit}

\textbf{Worum es ging:} Gliederung, Formulierung, Kürzung und LaTeX-Feinschliff (Tabellen, Abbildungen, Literatur, Overfull-Boxen).

\textbf{Beispielhafte Prompts:}
\begin{itemize}
  \item \enquote{Dieser Absatz über die Domänentrennung ist zu lang. Gib mir ein Beispiel, wie ich ihn straffen kann, ohne die Kernbegriffe COLMAP, Maskenwarp und ideale Domäne zu verlieren.}
  \item \enquote{Formuliere diesen Satz wissenschaftlicher und ohne Füllwörter: ,Das flexiblere Modell ist nicht grundsätzlich genauer, weil ...' -- gib zwei Alternativen.}
  \item \enquote{Die Tabelle \texttt{tab:bildablagen} ist zu breit. Formatiere die Tabelle neu, ohne Inhalt zu verlieren.}
\end{itemize}

\textbf{Prüfung:} Alle Textvorschläge wurden satzweise gelesen, fachlich geprüft und in eigenen Worten überarbeitet; Zitate und Zahlen wurden gegen Manifeste und Tabellen verifiziert.

\section{Fehleranalyse und Auswertung}

\textbf{Worum es ging:} Diagnose fehlgeschlagener Läufe (leere Masken, SUGAR-Importfehler, COLMAP-Registrierung), Einordnung von Screenshots und Zwischenständen.

\textbf{Beispielhafter Prompt:}
\begin{itemize}
  \item \enquote{Hier ein Log-Auszug mit leeren Eval-Masken. Welche Prüfung fehlt, und wo würdest du ein Quality-Gate einbauen? Erkläre den Vorschlag, bevor du Code änderst.}
\end{itemize}

\textbf{Prüfung:} Ursachen wurden im Code und in den Logs nachvollzogen; Korrekturen wurden erst nach manuellem Test übernommen.

\section{Grafiken und Auswertung}

\textbf{Worum es ging:} Erstellung der Panels (z.\,B.\ \texttt{pa\_panel\_sugar\_iteration}), Metrikplots (STS vs.\ SuGaR, Delta, Laufzeit) und deren konsistente Beschriftung.

\textbf{Beispielhafte Prompts:}
\begin{itemize}
  \item \enquote{Das Panel hat 9200er-Labels aus den historischen Dateinamen, soll aber konsistent 9001 zeigen. Schreibe ein Skript, das die Kacheln neu zusammensetzt und nur 9001 ausgibt.}
  \item \enquote{Lies aus den archivierten \texttt{sts\_masked.json}-Dateien die objektmaskierten Werte PSNR, SSIM und LPIPS mit einem neuen Skript aus und erzeuge daraus eine übersichtliche Grafik je Kameramodell.}
  \item \enquote{Prüfe alle Metrikgrafiken auf konsistente Farb- und Symbolkodierung und korrigiere sie über alle Grafiken hinweg: \texttt{OPENCV} immer orange, \texttt{SIMPLE\_RADIAL} immer blau, \texttt{PINHOLE} immer grau; je Auflösung ein eigener Formtyp (720p, QHD, low). Die Kodierung muss in jeder Grafik derselben Logik folgen.}
\end{itemize}

\textbf{Prüfung:} Generierte Grafiken wurden visuell und gegen die Quelldaten (JSON/CSVs) geprüft.

\section{Eigenständigkeit}

Alle Inhalte wurden fachlich geprüft, vollständig verstanden und können selbstständig vertreten werden. KI-Ausgaben wurden nicht als wissenschaftliche Quellen zitiert. Die Verantwortung für fachliche Richtigkeit und Eigenständigkeit liegt beim Verfasser.

\vfill
\noindent Würzburg, den \rule{4cm}{0.4pt}\hfill
\rule{6cm}{0.4pt}\\
\hfill Unterschrift des/der Studierenden

\end{document}
